% Generated from tex/chapters/ch04/05_ソフトウェア構築の基礎.tex for the SWEBOK structure tree
\subsubsection{複雑さの最小化}

人間の作業記憶には限界がある。複雑な構造や大量の詳細を長時間正確に保持することは難しい。この限界は、ソフトウェア構築の中心課題である。コードが複雑になると、理解しにくく、テストしにくく、変更しにくくなる。したがって、構築では「賢そうなコード」よりも「単純で読みやすいコード」を優先する必要がある。

複雑さには複数の種類がある。条件分岐やループが多いコードは、実行経路が多くなる。依存関係が多いコードは、変更の影響範囲が広くなる。名前がわかりにくいコードは、読者の理解を妨げる。ビルド手順が複雑であれば、構築そのものが不安定になる。

循環的複雑度は、コードの制御フローの複雑さを測る代表的な指標である。大まかには、独立した実行経路の数を表す。値が大きいほど、理解やテストが難しくなる傾向がある。ただし、指標は判断の補助であり、指標だけでよいコードかどうかを決めてはならない。

複雑さを下げる実務上の手段には、命名規則、適切な関数分割、モジュール化、重複の排除、標準的な制御構造の利用、読みやすいレイアウト、レビュー、静的解析、単体テストがある。

\begin{keybox}{見分け方}
複雑なコードは、説明に時間がかかる。単純なコードは、名前、構造、流れから意図が読み取れる。構築では「動くか」だけでなく「次の人が読めるか」を基準にする必要がある。
\end{keybox}
